\documentclass{article}
\usepackage{lmodern}
\usepackage{amsmath}
\usepackage{listings}
\usepackage{fullpage}
\usepackage{parskip}
\usepackage{xcolor}
\lstset{
	basicstyle=\ttfamily,
	columns=fullflexible,
	frame=single,
	breaklines=true,
	postbreak=\mbox{\textcolor{red}{$\hookrightarrow$}\space},
}
\begin{document}
\title{Experiment `p\_6\_1\_c\_process\_array\_reverse\_order\_array' Results}
\author{\today}
\date{}
\maketitle
\textbf{Experiment outcome: }SUCCESS
\\\textbf{Bad responses: }0
\\\textbf{Responses containing}~\texttt{assume}~\textbf{: }0
\\\textbf{Responses containing}~\texttt{expect}~\textbf{: }0
\\\textbf{Resolution attempts: }1
\\\textbf{Hard fails (resolution): }0
\\\textbf{Soft fails (resolution): }0
\\\textbf{Verification attempts: }1
\section*{Problem Specification}
\textbf{Problem name: }p\_6\_1\_c\_process\_array\_reverse\_order\_array
\\\textbf{Natural language statement: }Write a method that takes an array of ten integers and then returns all elements in reverse order.
\\\textbf{Method signature: }p\_6\_1\_c\_process\_array\_reverse\_order\_array(inputs: array<int>) returns (reversed: array<int>)
\subsection*{Ensures}
\begin{itemize}
\item \texttt{reversed.Length == inputs.Length}
\item \texttt{forall i :: 0 <= i < reversed.Length ==> reversed[i] == inputs[inputs.Length - 1 - i]}
\end{itemize}
\subsection*{Requires}
\begin{itemize}
\item \texttt{inputs.Length == 10}
\end{itemize}
\clearpage
\section*{GenAI interactions}
Below you will find all interactions between the `user' (program) and the `assistant' (OpenAI).
\subsection*{Program $\rightarrow$ GenAI}
\begin{lstlisting}
You are given the following task to perform in Dafny:

Write a method that takes an array of ten integers and then returns all elements in reverse order.

The signature should be:

p_6_1_c_process_array_reverse_order_array(inputs: array<int>) returns (reversed: array<int>)

The method should respect the following contract:

requires inputs.Length == 10, ensures reversed.Length == inputs.Length, ensures forall i :: 0 <= i < reversed.Length ==> reversed[i] == inputs[inputs.Length - 1 - i]

Produce and show only the Dafny body of this method, including the curly braces that surround it. Do not show the signature nor contract. You must not use 'assume' nor 'expect' anywhere in your code.
\end{lstlisting}
\subsection*{GenAI $\rightarrow$ Program}
\textbf{System fingerprint: }null
\\\textbf{ID: }chatcmpl-Dn09UE89fJkauMiJEnWHZlmIuZ6PB
\\\textbf{Model: }gpt-5.5-2026-04-23
\\\textbf{Created at: }1780571028
\\\textbf{Finish reason: }stop
\\\textbf{Usage: }344 tokens in, and 355 tokens out
\begin{lstlisting}
{
  reversed := new int[inputs.Length];
  var i := 0;
  while i < reversed.Length
    invariant reversed.Length == inputs.Length
    invariant 0 <= i <= reversed.Length
    invariant forall j :: 0 <= j < i ==> reversed[j] == inputs[inputs.Length - 1 - j]
  {
    reversed[i] := inputs[inputs.Length - 1 - i];
    i := i + 1;
  }
}
\end{lstlisting}
\section*{Final Program}
\begin{lstlisting}
method p_6_1_c_process_array_reverse_order_array(inputs: array<int>) returns (reversed: array<int>)
	requires inputs.Length == 10
	ensures reversed.Length == inputs.Length
	ensures forall i :: 0 <= i < reversed.Length ==> reversed[i] == inputs[inputs.Length - 1 - i]
{
  reversed := new int[inputs.Length];
  var i := 0;
  while i < reversed.Length
    invariant reversed.Length == inputs.Length
    invariant 0 <= i <= reversed.Length
    invariant forall j :: 0 <= j < i ==> reversed[j] == inputs[inputs.Length - 1 - j]
  {
    reversed[i] := inputs[inputs.Length - 1 - i];
    i := i + 1;
  }
}
\end{lstlisting}
\section*{Total Token Usage}
\textbf{Input tokens: }344
\\\textbf{Output tokens: }355
\\\textbf{Reasoning tokens: }231
\\\textbf{Sum of `total tokens': }699
\section*{Experiment Timings}
\textbf{Overall Experiment} started at 1780571020464, ended at 1780571062243, lasting 41779ms (41.78 seconds)
\\\textbf{Iteration \#1} started at 1780571020481, ended at 1780571062240, lasting 41759ms (41.76 seconds)
\end{document}
